\section{Materials and Methods}
\label{sec:methods}

\subsection{Dataset}
\label{sec:dataset}

The Cleveland subset of the UCI Heart Disease dataset~\cite{detrano1989cleveland} was used, comprising 242 patient records and 13 clinical attributes. The binary classification target indicates the presence (\texttt{num}~=~1) or absence (\texttt{num}~=~0) of coronary artery disease. The 13 features include demographic descriptors (age, sex), clinical measurements (resting blood pressure \texttt{trestbps}, serum cholesterol \texttt{chol}, fasting blood sugar \texttt{fbs}, maximum heart rate \texttt{thalach}, ST depression \texttt{oldpeak}), categorical clinical findings (chest pain type \texttt{cp}, resting ECG \texttt{restecg}, exercise-induced angina \texttt{exang}, ST-segment slope \texttt{slope}, thallium scan \texttt{thal}), and the number of major vessels affected on fluoroscopy (\texttt{ca}). Two missing values were present in the dataset --- one each in \texttt{ca} and \texttt{thal} --- and one extreme outlier was identified in \texttt{chol} (564~mg/dL). The class distribution is mildly imbalanced (54.1\% no-disease, 45.9\% disease), corresponding to an imbalance ratio of approximately 1.18:1.

\subsection{Preprocessing}
\label{sec:preprocessing}

Preprocessing was implemented as an \texttt{sklearn} \texttt{ColumnTransformer} that applies different transformations to feature subsets. Continuous features (\texttt{age}, \texttt{trestbps}, \texttt{chol}, \texttt{thalach}, \texttt{oldpeak}) were standardised to zero mean and unit variance via \texttt{StandardScaler}, with mean imputation applied beforehand as a defensive default. Categorical and binary features were imputed using the most-frequent (mode) value, addressing the two missing values in \texttt{ca} and \texttt{thal} without introducing distortion to the other categorical features. No scaling was applied to categorical features, as this would corrupt their integer encoding.

The transformer was constructed as a fresh, unfitted instance for each cross-validation fold and incorporated as the first step of a pipeline that ended with the classifier. This design ensures that imputation statistics and scaling parameters are computed exclusively on training data within each fold, with no leakage from validation or test data.

\subsection{Repeated Nested Cross-Validation}
\label{sec:rncv}

Generalisation performance was estimated via repeated nested cross-validation (rnCV) using $R=10$ independent repetitions of $N=5$-fold outer~$\times$~$K=3$-fold inner stratified cross-validation. This produces 50 outer-test scores per algorithm per metric. Stratified splitting preserved the class ratio in every fold. To ensure reproducibility, the random seed for repetition $r$ was set to $r_{\text{seed}} = r_{\text{base}} + r$, with $r_{\text{base}} = 42$, propagating identically to fold-splitters, classifier seeds, and the Optuna sampler. Figure~\ref{fig:ncv_diagram} illustrates one iteration of the pipeline.

\begin{figure}[t]
    \centering
    \includegraphics[width=\columnwidth]{figures/ncv_pipeline_diagram.pdf}
    \caption{One iteration of the rnCV pipeline. The outer loop estimates generalisation performance; the inner loop performs hyperparameter tuning. Preprocessing and (in Task~4) feature selection are fitted strictly on training folds at each level.}
    \label{fig:ncv_diagram}
\end{figure}

The inner loop performs hyperparameter optimisation via Optuna~\cite{akiba2019optuna} using the Tree-structured Parzen Estimator (TPE) sampler with 50 trials per inner study. For each candidate hyperparameter configuration, the inner-validation MCC is averaged across the three inner folds before being returned to the optimiser. Aggregating across folds before optimisation favours hyperparameters that generalise across inner-train splits over those that exploit a single favourable split. The Matthews Correlation Coefficient (MCC) was selected as the inner-loop metric for three reasons: (i) it incorporates all four cells of the confusion matrix, (ii) it is invariant under class swap, and (iii) it remains well-defined under mild class imbalance, unlike F1.

The complete metric battery computed on each outer-test fold comprised MCC, ROC-AUC, balanced accuracy, F1, recall, specificity, precision, and PR-AUC. For each algorithm and metric, the median across the 50 outer-test scores was reported as the central estimate, with non-parametric 95\% confidence intervals (CIs) constructed by percentile bootstrap resampling (10,000 resamples). Pairs of algorithms with overlapping CIs on median MCC were treated as statistically indistinguishable on this primary metric.

\subsection{Algorithm Comparison}
\label{sec:algocomp}

Seven classifiers spanning linear, probabilistic, discriminant, and ensemble paradigms were compared: Logistic Regression with Elastic Net regularisation (LR), Gaussian Naive Bayes (GNB), Linear Discriminant Analysis with shrinkage (LDA), Random Forest (RF), LightGBM, XGBoost, and CatBoost. Each algorithm was first evaluated under default hyperparameters using $R=10$ iterations of 5-fold stratified CV (the inner loop disabled, $K = \emptyset$), establishing a fair baseline against which the gain from tuning could be assessed. The full rnCV pipeline was then run with algorithm-specific hyperparameter spaces summarised in Table~\ref{tab:hp_spaces}. Ranges for regularisation parameters were defined on a log scale to reflect their multiplicative nature; integer ranges (tree depth, number of estimators) covered values that are conventional for small tabular datasets.

\begin{table}[t]
    \centering
    \caption{Hyperparameter search spaces per algorithm.}
    \label{tab:hp_spaces}
    \resizebox{\columnwidth}{!}{%
    \begin{tabular}{lll}
        \toprule
        \textbf{Algorithm} & \textbf{Hyperparameter} & \textbf{Range} \\
        \midrule
        LR  & C            & log $[10^{-3}, 10^{1}]$ \\
            & l1\_ratio    & $[0, 1]$ \\
        \midrule
        GNB & var\_smoothing & log $[10^{-12}, 10^{-6}]$ \\
        \midrule
        LDA & solver       & \{lsqr, eigen\} \\
            & shrinkage    & $[0, 1]$ \\
        \midrule
        RF  & n\_estimators & $\{100,\ldots,500\}$ step 50 \\
            & max\_depth   & $\{3,\ldots,20\}$ \\
            & min\_samples\_split & $\{2,\ldots,20\}$ \\
            & min\_samples\_leaf  & $\{1,\ldots,10\}$ \\
            & max\_features      & \{sqrt, log2\} \\
        \midrule
        LightGBM & n\_estimators & $\{100,\ldots,500\}$ step 50 \\
            & learning\_rate    & log $[10^{-3}, 0.3]$ \\
            & num\_leaves       & $\{8,\ldots,64\}$ \\
            & max\_depth        & $\{3,\ldots,12\}$ \\
            & min\_child\_samples & $\{5,\ldots,30\}$ \\
            & reg\_alpha, reg\_lambda & log $[10^{-3}, 1]$ \\
        \midrule
        XGBoost & n\_estimators & $\{100,\ldots,500\}$ step 50 \\
            & learning\_rate     & log $[10^{-3}, 0.3]$ \\
            & max\_depth         & $\{3,\ldots,10\}$ \\
            & min\_child\_weight & $\{1,\ldots,10\}$ \\
            & subsample, colsample\_bytree & $[0.6, 1.0]$ \\
            & reg\_alpha, reg\_lambda & log $[10^{-3}, 1]$ \\
        \midrule
        CatBoost & iterations    & $\{100,\ldots,500\}$ step 50 \\
            & learning\_rate     & log $[10^{-3}, 0.3]$ \\
            & depth              & $\{3,\ldots,10\}$ \\
            & l2\_leaf\_reg      & $[1, 10]$ \\
        \bottomrule
    \end{tabular}%
    }
\end{table}

The winning algorithm was selected by ranking on median MCC. Ties (overlapping 95\% CIs) were to be broken on median AUC and PR-AUC.

\subsection{Feature Selection}
\label{sec:feature_selection}

Following identification of the winning algorithm, model-agnostic feature selection was performed using permutation importance~\cite{breiman2001random,fisher2019all} with MCC as the scoring metric. The selector was applied within the outer loop of the rnCV: for each outer training fold, a default-hyperparameter LDA was fitted, permutation importances were estimated with 10 repeats per feature, and the top-$k$ features were retained. The same feature subset was then propagated to the inner loop and final outer-test evaluation. The locked outer test fold was never seen by the selector at any stage.

The number of retained features $k$ was determined by sweeping $k \in \{3, 4, \ldots, 12\}$ and selecting the value maximising median MCC across the 50 outer-test scores. This data-driven optimum was preferred over a more parsimonious criterion (e.g., the smallest $k$ whose 95\% CI overlaps the full-feature CI) because it captures the most predictive feature set the data supports without over-parametrisation. To assess feature stability, the selection frequency of each feature was recorded across all 50 fold $\times$ repetition outcomes at the chosen $k$. Features selected in $\geq 80\%$ of folds, in line with conventions from the stability selection literature~\cite{meinshausen2010stability}, were retained as the \emph{stable feature set} for downstream deployment.

\subsection{Final Model and Persistence}
\label{sec:final_model}

The final model was trained on the entire dataset using the stable feature set. Optimal hyperparameters were determined by a single (non-nested) 5-fold stratified CV over the LDA hyperparameter space using Optuna with 50 trials. The deployable pipeline was constructed in three stages: (i) a custom \texttt{FeatureSubsetSelector} accepting raw 13-column input and projecting onto the stable feature subset, (ii) the \texttt{ColumnTransformer} for imputation and standardisation, and (iii) the LDA classifier with the chosen hyperparameters. The complete pipeline was fitted on all 242 patients and serialised using \texttt{pickle} as a single artefact (\texttt{models/final\_model.pkl}). To verify deployability, the pickled artefact was reloaded in a fresh state and verified to produce predictions on raw 13-column input without external preprocessing.

\subsection{Model Interpretation with SHAP}
\label{sec:shap}

SHAP (SHapley Additive exPlanations) values~\cite{lundberg2017unified} were computed on the final model fitted on all available data, providing model interpretation grounded in cooperative game theory. The model-agnostic \texttt{KernelExplainer} was selected, as LDA is neither tree-based (precluding \texttt{TreeExplainer}) nor straightforwardly compatible with SHAP's \texttt{LinearExplainer} due to its probabilistic output structure. To keep computation tractable, the background distribution was summarised using $k$-means clustering with 50 centroids; this approximation preserves the input distribution while reducing the cost of the marginal-expectation estimates.

SHAP values were computed for each of the 242 patients on the stable feature subset, yielding (i) a global ranking of mean absolute SHAP per feature, (ii) a beeswarm summary plot showing both magnitude and direction of effect, and (iii) waterfall plots decomposing individual predictions into per-feature contributions for two illustrative cases (one correctly predicted disease, one correctly predicted healthy).

\subsection{Bonus Task: Error Analysis}
\label{sec:bonus_methods}

A complementary error analysis was performed using a single 70/30 stratified train/validation split (\texttt{random\_state}~=~42), with the model architecture and hyperparameters fixed to those of the final model. After predicting on the validation set, each patient was assigned to one of four categories: true positive (TP), true negative (TN), false positive (FP), or false negative (FN). Demographic and clinical features were compared across these categories using Mann-Whitney U tests for continuous variables and Fisher's exact tests for categorical variables. SHAP values were computed on the validation set, and per-feature SHAP distributions were compared across the four categories to identify features whose contribution patterns differ between correct and incorrect predictions. Subgroup-specific error rates were computed by stratifying on sex, age band, and key clinical features.

\subsection{Implementation and Reproducibility}
\label{sec:implementation}

The pipeline was implemented in Python 3.10 using \texttt{scikit-learn} (1.3), \texttt{optuna} (3.4), \texttt{shap} (0.44), \texttt{lightgbm}, \texttt{xgboost}, and \texttt{catboost}. The code follows object-oriented design, with the rnCV procedure encapsulated in a single \texttt{RepeatedNestedCV} class (in \texttt{src/rncv.py}) that accepts a list of \texttt{EstimatorConfig} objects pairing each estimator with its hyperparameter-space callable. Random seeds were fixed throughout. \texttt{Optuna} was set to single-thread execution to avoid stochasticity in the TPE sampler; per-tree parallelism (\texttt{n\_jobs}~=~$-1$) was retained only for Random Forest, where each tree is independently seeded. The full repository is available at: \href{https://github.com/GiorgosBoulogeorgos/Machine-learning-in-computational-biology-assignment2-}{GitHub Repository}.
